\documentclass[11pt]{article}
\usepackage[margin=1in]{geometry}
\usepackage{amsmath,amssymb,amsthm,mathtools}
\usepackage[colorlinks=true,linkcolor=blue,citecolor=blue]{hyperref}

\newtheorem{theorem}{Theorem}
\newtheorem{lemma}{Lemma}
\newtheorem{corollary}{Corollary}
\theoremstyle{definition}
\newtheorem{assumption}{Standing fact}
\theoremstyle{remark}
\newtheorem*{remark}{Remark}

\newcommand{\dd}{\,\mathrm{d}}
\newcommand{\Ynode}{Y_{\mathrm{node}}}
\newcommand{\TW}{\mathrm{TW}}
\newcommand{\R}{\mathbb{R}}
\newcommand{\eqd}{\stackrel{d}{=}}
\newcommand{\xit}{\widetilde{\xi}}
\DeclareMathOperator{\Ai}{Ai}

\title{\textbf{A complete proof of the shooting characterisation:\\
the noisy-Airy inner exit measure is Tracy--Widom}}
\author{Solomon Williams \\ \small University of Edinburgh}
\date{\today}

\begin{document}
\maketitle

\begin{abstract}
\noindent
We prove that the inner exit measure of the noisy folded limit cycle---the law of the
first-node level $\Ynode$ of the swept noisy Airy equation---coincides with the law of
minus the ground-state eigenvalue of the stochastic Airy operator, and hence with
Tracy--Widom: $\Ynode\eqd-\Lambda_0(\beta)\eqd\TW_\beta$, $\eta=2/\sqrt\beta$. The proof
is elementary given the random-Schr\"odinger framework of Ram\'irez--Rider--Vir\'ag
(RRV): the swept solution is the \emph{energy-zero} solution of a random Schr\"odinger
operator $M$, its first zero is the level at which the truncated $M$ has ground state
$0$, and that level is read off through a \emph{monotone spectral function} whose
defining event involves only a \emph{deterministic} translation---so the stationarity
and reflection symmetry of white noise close the identity. The single nontrivial random
shift, which obstructed earlier arguments, is bypassed by Dirichlet domain
monotonicity.
\end{abstract}

\section{Setup and standing facts}

\paragraph{The two objects.} Fix $\eta>0$ and $\beta=4/\eta^2$. Let $\xit$ be a white
noise on $\R$ (the distributional derivative of a two-sided Brownian motion $W$).

\emph{Object 1 (the inner exit level).} In the rescaling chart the noisy canard solves,
in the slow variable $Y$,
\begin{equation}\label{eq:Yeq}
  u_{YY}=\bigl(Y-\eta\,\xit\bigr)\,u .
\end{equation}
(This is \eqref{eq:Yeq} because $Y=Y_0-T$ gives $\partial_Y^2=\partial_T^2$, and the
time white noise becomes a white noise in $Y$.) Let $\psi$ be the solution of
\eqref{eq:Yeq} that is recessive (square-integrable) as $Y\to+\infty$; its largest zero
is the \emph{inner exit level} $\Ynode$.

\emph{Object 2 (the SAO ground state).} The stochastic Airy operator
\begin{equation}\label{eq:SAO}
  H_\beta=-\frac{\dd^2}{\dd x^2}+x+\frac{2}{\sqrt\beta}\,b'(x)
  \qquad\text{on }[0,\infty),\ \text{Dirichlet at }0,
\end{equation}
($b'$ a white noise, $2/\sqrt\beta=\eta$) has a.s.\ a simple lowest eigenvalue
$\Lambda_0(\beta)$.

\paragraph{The operator $M$.} Write \eqref{eq:Yeq} as $Mu=0$ with the random
Schr\"odinger operator
\begin{equation}\label{eq:M}
  M=-\frac{\dd^2}{\dd Y^2}+Y-\eta\,\xit(Y),
\end{equation}
i.e.\ $\psi$ is the energy-$0$ recessive solution of $M$. For $\ell\in\R$ let
$M_\ell$ denote $M$ acting on $L^2([\ell,\infty))$ with a Dirichlet condition at $\ell$.

We use four standing facts, each classical in its area.

\begin{assumption}[white-noise invariances]\label{F1}
For every deterministic $c\in\R$, $\ \xit(c+\cdot)\eqd\xit\ $ (stationarity) and
$-\xit\eqd\xit$ (reflection symmetry). Consequently $-\eta\,\xit(c+\cdot)\eqd
\tfrac{2}{\sqrt\beta}b'$.
\end{assumption}

\begin{assumption}[random Schr\"odinger operator]\label{F2}
For a.e.\ realisation and every $\ell$, $M_\ell$ is self-adjoint and bounded below with
purely discrete spectrum near its bottom; its lowest eigenvalue is simple with a
ground-state eigenfunction that is strictly positive on $(\ell,\infty)$. The same holds
for $H_\beta$. \emph{(This is the RRV construction \cite{RRV}; the confining potential
$Y\to+\infty$ and the white-noise term are handled by their quadratic-form/Riccati
analysis. See also \cite{FN} for $1$D Schr\"odinger operators with white-noise potential.)}
\end{assumption}

\begin{assumption}[Weyl limit point at $+\infty$]\label{F3}
Since the potential $\to+\infty$, equation \eqref{eq:Yeq} is in the limit-point case at
$+\infty$: for each energy there is a unique-up-to-scale solution that is square-integrable
near $+\infty$ \cite{CL}. In particular the recessive solution $\psi$ exists and is
unique up to scale.
\end{assumption}

\begin{assumption}[RRV edge]\label{F4}
$-\Lambda_0(\beta)\eqd\TW_\beta$ \cite{RRV}.
\end{assumption}

\section{The theorem}

\begin{theorem}[Shooting characterisation]\label{thm:main}
With the notation above,
\[
  \Ynode \eqd -\Lambda_0(\beta)\eqd \TW_\beta ,\qquad \eta=2/\sqrt\beta .
\]
Moreover, if $u^{Y_0}$ denotes the solution of \eqref{eq:Yeq} started at $Y=Y_0$ on the
canard branch ($u_Y/u=-\sqrt{Y_0}$) and $\Ynode(Y_0)$ its largest zero, then
$\Ynode(Y_0)\to\Ynode$ almost surely as $Y_0\to\infty$; hence
$\Ynode(Y_0)\xrightarrow{d}-\Lambda_0(\beta)$.
\end{theorem}

The proof is four lemmas. The crux is Lemma~\ref{lem:mono}: the event $\{\Ynode\le t\}$
equals $\{G(t)\ge0\}$ for a monotone spectral function $G$, and the level $t$ there is
\emph{deterministic}, so Fact~\ref{F1} applies.

\begin{lemma}[the exit level is well defined]\label{lem:exists}
Almost surely $\psi$ exists (Fact~\ref{F3}), is real and not identically zero, is
positive for all large $Y$, and has a largest zero $\Ynode\in\R$.
\end{lemma}

\begin{proof}
Existence and uniqueness up to scale are Fact~\ref{F3}; take $\psi$ real (the equation
has real coefficients). For $Y$ large the potential $Y-\eta\xit$ is positive, so by
\eqref{eq:Yeq} $\psi$ is convex where positive and concave where negative; the recessive
(decaying) branch is positive for all large $Y$. For $Y\to-\infty$ the potential
$\to-\infty$ and $\psi$ oscillates with infinitely many zeros, which are isolated (a
nontrivial solution of a second-order linear ODE has only simple, isolated zeros). The
set of zeros is therefore bounded above (none for large $Y$) and discrete, so it has a
maximum $\Ynode$.
\end{proof}

\begin{lemma}[forgetting: the sweep limit]\label{lem:forget}
Almost surely $\Ynode(Y_0)\to\Ynode$ as $Y_0\to\infty$, exponentially fast.
\end{lemma}

\begin{proof}
Work with the Riccati variable $P=u_Y/u$, which solves
$P_Y=(Y-\eta\xit)-P^2$. On the barrier $\{Y>0\}$ the two WKB branches are
$P\approx\pm\sqrt{Y}$, and linearising about $P=-\sqrt Y$ gives
$\partial_P\bigl[(Y-\eta\xit)-P^2\bigr]=-2P\approx+2\sqrt Y$. Thus along \emph{decreasing}
$Y$ the branch $P=-\sqrt Y$ is attracting: if $P^{(1)},P^{(2)}$ are two solutions with
$\delta=P^{(1)}-P^{(2)}$ then $\delta_Y=-(P^{(1)}+P^{(2)})\delta$, so for $Y$ decreasing
from $Y_0$,
\[
  |\delta(Y)|=|\delta(Y_0)|\exp\!\Bigl(-\!\int_Y^{Y_0}\!(P^{(1)}+P^{(2)})\dd Y'\Bigr)
            \le|\delta(Y_0)|\,e^{-2\int_Y^{Y_0}\sqrt{Y'}\,(1+o(1))\dd Y'} .
\]
The canard datum $P^{Y_0}(Y_0)=-\sqrt{Y_0}$ differs from the recessive value
$P_\psi(Y_0)$ by the WKB error $\delta(Y_0)=O(Y_0^{-1})$ (Fact~\ref{F3}); integrating
down to any fixed $Y_*$ contracts this to $O\!\bigl(Y_0^{-1}e^{-\frac43 Y_0^{3/2}}\bigr)$,
uniformly. Hence $P^{Y_0}\to P_\psi$ locally uniformly on $\{Y\le Y_*\}$, and the largest
zero $\Ynode(Y_0)$ (a continuous functional of $P$ on a neighbourhood of $\Ynode$, where
$P\to-\infty$ transversally) converges to $\Ynode$.
\end{proof}

\begin{lemma}[monotone spectral function]\label{lem:mono}
Let $G(\ell)$ be the lowest eigenvalue of $M_\ell$ (Fact~\ref{F2}). Almost surely:
\begin{enumerate}
\item[\textup{(i)}] $G$ is strictly increasing on $\R$;
\item[\textup{(ii)}] $G(\Ynode)=0$;
\item[\textup{(iii)}] for every $t\in\R$,\quad $\{\Ynode\le t\}=\{G(t)\ge0\}$.
\end{enumerate}
\end{lemma}

\begin{proof}
\textit{(i)} For $\ell_1<\ell_2$, the form domain of $M_{\ell_2}$ embeds into that of
$M_{\ell_1}$ by extension by zero, and the quadratic forms agree on the image; hence by
the min--max principle $G(\ell_1)\le G(\ell_2)$ (Dirichlet domain monotonicity
\cite[XIII]{RS}). The inequality is strict: the ground state $\varphi_1$ of $M_{\ell_1}$
is strictly positive on $(\ell_1,\infty)$ (Fact~\ref{F2}), so $\varphi_1(\ell_2)>0$ and
$\varphi_1$ is \emph{not} admissible for $M_{\ell_2}$ (which requires vanishing at
$\ell_2$); the variational minimum therefore strictly increases.

\textit{(ii)} Restrict $\psi$ to $[\Ynode,\infty)$. There $\psi(\Ynode)=0$ (Dirichlet
satisfied), $\psi\in L^2$ (recessive decay at $+\infty$), $M\psi=0$, and $\psi$ has no
zero in $(\Ynode,\infty)$ since $\Ynode$ is the \emph{largest} zero
(Lemma~\ref{lem:exists}). A nodeless $L^2$ eigenfunction is the ground state
(Fact~\ref{F2}), so $0$ is the lowest eigenvalue of $M_{\Ynode}$: $G(\Ynode)=0$.

\textit{(iii)} By (i) and (ii), $G(t)\ge0=G(\Ynode)\iff t\ge\Ynode$.
\end{proof}

\begin{lemma}[deterministic translation]\label{lem:translate}
For each fixed $t\in\R$,\quad $G(t)\eqd\Lambda_0(\beta)+t$.
\end{lemma}

\begin{proof}
Fix $t$ and substitute the \emph{deterministic} shift $Y=t+x$, $x\ge0$. Then
\[
  M_t=-\frac{\dd^2}{\dd x^2}+(t+x)-\eta\,\xit(t+x)
     =\underbrace{\Bigl[-\frac{\dd^2}{\dd x^2}+x-\eta\,\xit(t+x)\Bigr]}_{=:H'_t}+\,t ,
\]
an operator on $[0,\infty)$ with Dirichlet at $0$. By Fact~\ref{F1}, $\xit(t+\cdot)\eqd
\xit$ and $-\eta\,\xit(t+\cdot)\eqd\tfrac{2}{\sqrt\beta}b'$, so $H'_t\eqd H_\beta$ as
random operators; in particular their lowest eigenvalues have the same law,
$\Lambda_0(H'_t)\eqd\Lambda_0(\beta)$. Since $M_t=H'_t+t$, $\;G(t)=\Lambda_0(H'_t)+t$,
whence $G(t)\eqd\Lambda_0(\beta)+t$.
\end{proof}

\begin{proof}[Proof of Theorem~\ref{thm:main}]
Fix $t\in\R$. Using Lemma~\ref{lem:mono}(iii) (a pathwise event identity) and then
Lemma~\ref{lem:translate} (a statement about the law of $G(t)$ at the fixed level $t$),
\[
  \mathbb P(\Ynode\le t)
  =\mathbb P\bigl(G(t)\ge0\bigr)
  =\mathbb P\bigl(\Lambda_0(\beta)+t\ge0\bigr)
  =\mathbb P\bigl(\Lambda_0(\beta)\ge-t\bigr)
  =\mathbb P\bigl(-\Lambda_0(\beta)\le t\bigr).
\]
As $t$ was arbitrary, $\Ynode\eqd-\Lambda_0(\beta)$; by Fact~\ref{F4},
$\Ynode\eqd\TW_\beta$. The sweep statement is Lemma~\ref{lem:forget} together with the
continuous-mapping theorem.
\end{proof}

\section{Discussion}

\paragraph{What is elementary and what is cited.} Given Facts~\ref{F1}--\ref{F4}, the
proof is elementary: Lemmas~\ref{lem:exists}, \ref{lem:mono} and \ref{lem:translate} use
only Sturm theory, Dirichlet domain monotonicity, and the substitution $Y=t+x$;
Lemma~\ref{lem:forget} is a Gronwall estimate on the Riccati. The only heavy input is
Fact~\ref{F2}---the rigorous meaning of a Schr\"odinger operator with a white-noise
potential and the simplicity/positivity of its ground state. That is precisely the RRV
framework we are matching to, applied to $M$, which is an SAO \emph{in law}
(Lemma~\ref{lem:translate} at $t=0$); so the proof imports no machinery beyond what the
statement already presupposes.

\paragraph{Why earlier arguments stalled.} A direct ``reflection'' argument tries to
translate by the random level $\Ynode$, producing a noise $\xit(\Ynode+\cdot)$ shifted by
a noise-dependent amount---not a clean white noise. Lemma~\ref{lem:mono} removes this
obstruction: the event $\{\Ynode\le t\}$ is recast as $\{G(t)\ge0\}$ with $t$
\emph{deterministic}, after which Fact~\ref{F1} applies verbatim.

\paragraph{Numerical confirmation.} (i) The conclusion $\Ynode\eqd-\Lambda_0(\beta)$ is
confirmed by computing $-\Lambda_0$ as the smallest eigenvalue of the tridiagonal
discretisation of $H_\beta$ and $\Ynode$ by ODE shooting---two methods with no shared
code path---which agree in mean, variance and skewness at $\beta=1,2,4$, and match the
exact $\TW_\beta$ densities. (ii) The mechanism of Lemma~\ref{lem:mono} is confirmed
directly: for fixed realisations, $G(\ell)$ (smallest eigenvalue of the discretised
$M_\ell$) is monotone increasing and crosses $0$ at the recessive solution's largest
zero $\Ynode$ computed from the same noise.

\paragraph{Consequence.} The inner exit measure of the noisy folded limit cycle is
$\TW_\beta$ with $\beta=4/\eta^2$, exactly; the early-escape probability is the tail mass
$1-F_\beta(0)$, and the folded-limit-cycle canard escape thereby lies in the
Tracy--Widom / KPZ edge universality class.

\begin{thebibliography}{9}
\bibitem{RRV} J.~Ram\'irez, B.~Rider, B.~Vir\'ag, \emph{Beta ensembles, stochastic Airy
spectrum, and a diffusion}, J.\ Amer.\ Math.\ Soc.\ \textbf{24} (2011) 919--944.
\bibitem{FN} M.~Fukushima, S.~Nakao, \emph{On spectra of the Schr\"odinger operator with
a white Gaussian noise potential}, Z.\ Wahrsch.\ \textbf{37} (1977) 267--274.
\bibitem{CL} E.~Coddington, N.~Levinson, \emph{Theory of Ordinary Differential
Equations}, McGraw--Hill (1955), Ch.~9 (limit-point/limit-circle).
\bibitem{RS} M.~Reed, B.~Simon, \emph{Methods of Modern Mathematical Physics IV},
Academic Press (1978), \S XIII (min--max, Dirichlet--Neumann bracketing).
\bibitem{TW} C.~Tracy, H.~Widom, \emph{Level-spacing distributions and the Airy kernel},
Comm.\ Math.\ Phys.\ \textbf{159} (1994) 151--174.
\end{thebibliography}

\end{document}
